% GENERATED FILE. Edit canonical lessons, then run npm run generate:books.
% canonical-source-hash: fnv1a64:3436aa7da1cf2cdc
% canonical-lessons: FR-C176-arreter, FR-C176-changer, FR-C176-deplacer, FR-C176-enlever, FR-C176-tenir

\chapter{To stop, To change, To move, To take off, To hold}
\label{ch:fr-to-stop-to-change-to-move-to-take-off-to-hold}
\hlchaptermodality{176}

\begin{chapteropening}
\textbf{By the end of this chapter:} \emph{I can say to stop, to change, to move, to take off and to hold.}

Complete the last lesson of chapter 176: I can say to stop, to change, to move, to take off and to hold.
\end{chapteropening}

\section[arrêter]{arrêter — to stop}

\label{lesson:FR-C176-arreter}



\begin{quote}
Before the new one: say the French for to invite, then the French for to park.
\end{quote}

\subsection*{You'll want to know: arrêter}
\textbf{arrêter} — \textquotedblleft{}to stop\textquotedblright{}.

\begin{cousinweb}
\textbf{arrêter} $\leftarrow$ Old French \emph{arester} $\leftarrow$ Vulgar Latin \emph{*arrestare}, from \emph{ad} and \emph{restare}, \textquotedblleft{}to stay back\textquotedblright{} — inherited. (EtymDB 2.1)
\end{cousinweb}

\begin{grammarlens}[title={What you've built}]
One more word for this chapter.
\end{grammarlens}

\subsection*{Your turn}
\begin{itemize}
\raggedright
  \item \emph{Say it:} \emph{arrêter}
  \item \emph{Say it:} \emph{arrêter}, once more
  \item \emph{Say it:} say \emph{arrêter}
  \item \emph{Recall:} say the French for to invite, then the French for to park, then say \emph{arrêter} again
\end{itemize}

\begin{tcolorbox}[breakable,colback=teal!4,colframe=teal!35!black,title={Before you move on}]
What does arrêter mean? (\textquotedblleft{}To stop\textquotedblright{}.) Say it once more.
\end{tcolorbox}
\section[changer]{changer — to change}

\label{lesson:FR-C176-changer}



\begin{quote}
Before the new one: say the French for to park, then the French for to stop.
\end{quote}

\subsection*{You'll want to know: changer}
\textbf{changer} — \textquotedblleft{}to change\textquotedblright{}.

\begin{cousinweb}
\textbf{changer} $\leftarrow$ Old French \emph{changier} $\leftarrow$ Late Latin \emph{cambiare}, \textquotedblleft{}to exchange\textquotedblright{}, a word Latin took from Gaulish — inherited. (EtymDB 2.1)
\end{cousinweb}

\begin{grammarlens}[title={What you've built}]
One more word for this chapter.
\end{grammarlens}

\subsection*{Your turn}
\begin{itemize}
\raggedright
  \item \emph{Say it:} \emph{changer}
  \item \emph{Say it:} \emph{changer}, once more
  \item \emph{Say it:} say \emph{changer}
  \item \emph{Recall:} say the French for to park, then the French for to stop, then say \emph{changer} again
\end{itemize}

\begin{tcolorbox}[breakable,colback=teal!4,colframe=teal!35!black,title={Before you move on}]
What does changer mean? (\textquotedblleft{}To change\textquotedblright{}.) Say it once more.
\end{tcolorbox}
\section[déplacer]{déplacer — to move (something)}

\label{lesson:FR-C176-deplacer}



\begin{quote}
Before the new one: say the French for to stop, then the French for to change.
\end{quote}

\subsection*{You'll want to know: déplacer}
\textbf{déplacer} — \textquotedblleft{}to move (something)\textquotedblright{}.

\begin{cousinweb}
\textbf{déplacer} — built from \emph{dé-} and the word for \textquotedblleft{}a place\textquotedblright{}, $\leftarrow$ Latin \emph{platea}, \textquotedblleft{}a broad street\textquotedblright{} (Lewis \& Short).
\end{cousinweb}

\begin{grammarlens}[title={What you've built}]
One more word for this chapter.
\end{grammarlens}

\subsection*{Your turn}
\begin{itemize}
\raggedright
  \item \emph{Say it:} \emph{déplacer}
  \item \emph{Say it:} \emph{déplacer}, once more
  \item \emph{Say it:} say \emph{déplacer}
  \item \emph{Recall:} say the French for to stop, then the French for to change, then say \emph{déplacer} again
\end{itemize}

\begin{tcolorbox}[breakable,colback=teal!4,colframe=teal!35!black,title={Before you move on}]
What does déplacer mean? (\textquotedblleft{}To move (something)\textquotedblright{}.) Say it once more.
\end{tcolorbox}
\section[enlever]{enlever — to take off, to remove}

\label{lesson:FR-C176-enlever}



\begin{quote}
Before the new one: say the French for to change, then the French for to move.
\end{quote}

\subsection*{You'll want to know: enlever}
\textbf{enlever} — \textquotedblleft{}to take off, to remove\textquotedblright{}.

\begin{cousinweb}
\textbf{enlever} — built from \emph{en}, \textquotedblleft{}away\textquotedblright{}, and \emph{lever}, \textquotedblleft{}to lift\textquotedblright{}, $\leftarrow$ Latin \emph{levare}.
\end{cousinweb}

\begin{grammarlens}[title={What you've built}]
One more word for this chapter.
\end{grammarlens}

\subsection*{Your turn}
\begin{itemize}
\raggedright
  \item \emph{Say it:} \emph{enlever}
  \item \emph{Say it:} \emph{enlever}, once more
  \item \emph{Say it:} say \emph{enlever}
  \item \emph{Recall:} say the French for to change, then the French for to move, then say \emph{enlever} again
\end{itemize}

\begin{tcolorbox}[breakable,colback=teal!4,colframe=teal!35!black,title={Before you move on}]
What does enlever mean? (\textquotedblleft{}To take off, to remove\textquotedblright{}.) Say it once more.
\end{tcolorbox}
\section[tenir]{tenir — to hold}

\label{lesson:FR-C176-tenir}



\begin{quote}
Before the new one: say the French for to move, then the French for to take off.
\end{quote}

\subsection*{You'll want to know: tenir}
\textbf{tenir} — \textquotedblleft{}to hold\textquotedblright{}.

\begin{cousinweb}
\textbf{tenir} $\leftarrow$ Vulgar Latin \emph{*tenire} $\leftarrow$ Latin \emph{tenere}, \textquotedblleft{}to hold\textquotedblright{} — inherited. (EtymDB 2.1)
\end{cousinweb}

\begin{grammarlens}[title={What you've built}]
One more word for this chapter.
\end{grammarlens}

\subsection*{Your turn}
\begin{itemize}
\raggedright
  \item \emph{Say it:} \emph{tenir}
  \item \emph{Say it:} \emph{tenir}, once more
  \item \emph{Say it:} say \emph{tenir}
  \item \emph{Recall:} say the French for to move, then the French for to take off, then say \emph{tenir} again
\end{itemize}

\begin{tcolorbox}[breakable,colback=teal!4,colframe=teal!35!black,title={Before you move on}]
What does tenir mean? (\textquotedblleft{}To hold\textquotedblright{}.) Say it once more.
\end{tcolorbox}
